Teach a robot a manipulation skill from human demonstrations and you face a quiet
decision with loud consequences. Training writes many checkpoints to disk. Only one gets
deployed. Which one? Running every checkpoint on the robot is too expensive, so the field
defaults to a cheap proxy: pick the checkpoint with the lowest validation loss on held-out
expert actions. In image classification that proxy is near-perfect. In imitation learning
it is not.

The mismatch is not a secret. Robomimic~\citep{mandlekar2021} reported that the
validation-loss-best checkpoint can be 50 to 100 percent worse than the best checkpoint you
would have found with rollouts. ACT~\citep{zhao2023act} adopted validation-loss selection as a
working compromise, not a principled choice. SIMPLER~\citep{li2024simpler} measured low correlation
between validation error and real success. NeME~\citep{tiezzi2025neme} built a learned
alternative for one task family. Each group noticed the gap in passing. None characterised it.

That leaves three open questions. How big is the gap, and what makes it grow? Does any
offline quantity, computable without rollouts, track success better than validation loss?
And can a combination of such quantities generalise to tasks it was never tuned on? We
answer all three with a pre-registered study.

Here is the framing that organises the work. Validation loss is a surrogate endpoint;
success rate is the true endpoint. Calibrating one against the other is exactly the exercise
drug developers run when they ask whether LDL cholesterol predicts cardiovascular
mortality~\citep{prentice1989}. The surrogate is cheap, the true endpoint is expensive, and
the rigorous practice is to map one to the other across many situations. We import that
discipline into robot learning.

Our protocol is locked before any training: the tasks, the eight offline metrics, the
rollout count, the held-out split, and the analysis code are fixed in advance and hashed.
That matters because the result is publication-worthy whichever way it falls. If the gap is
large and regime-dependent, the failure boundary is a finding. If an offline metric recovers
the signal, training pipelines get cheaper. If none does, we have a calibrated negative
result that justifies the cost of rollout-based evaluation. The outcome cannot be an
uninterpretable null.
